\documentclass[11pt]{article}
\newcommand{\LabNumber}{3}
\newcommand{\LabTitle}{AVR I/O and Software Time Delay}
\newcommand{\LabTopic}{Bare-metal AVR assembly in Microchip Studio, hex upload to an Arduino Uno R3 (ATmega328P) via AVRdude, and oscilloscope-verified software time delays}
% Shared preamble for Computer Programming lab sheets.
% Include with % Shared preamble for Computer Programming lab sheets.
% Include with % Shared preamble for Computer Programming lab sheets.
% Include with \input{common.tex} after \documentclass{article}.

\usepackage[a4paper,margin=2.2cm]{geometry}
\usepackage[T1]{fontenc}
\usepackage{lmodern}
\usepackage{microtype}
\usepackage{enumitem}
\usepackage{xcolor}
\usepackage{tcolorbox}
\tcbuselibrary{breakable,skins}
\usepackage{listings}
\usepackage{fancyhdr}
\usepackage{titlesec}
\usepackage{amsmath}
\usepackage{array}
\usepackage{booktabs}
\usepackage{hyperref}

\hypersetup{
  colorlinks=true,
  linkcolor=black,
  urlcolor=blue!60!black,
  pdfborder={0 0 0}
}

% ---------- Course identity (override per file if needed) ----------
\providecommand{\CourseName}{Microprocessors}
\providecommand{\CourseCode}{010153521}
\providecommand{\Semester}{Semester 1/2026}
\providecommand{\LabDuration}{3 hours (in-lab)}

% ---------- Colors ----------
\definecolor{accent}{HTML}{1F4E79}
\definecolor{accentlight}{HTML}{D9E2EC}
\definecolor{warn}{HTML}{B45309}
\definecolor{ok}{HTML}{166534}
\definecolor{codebg}{HTML}{F6F8FA}
\definecolor{codeframe}{HTML}{D0D7DE}
\definecolor{codekw}{HTML}{0550AE}
\definecolor{codestr}{HTML}{0A7B2C}
\definecolor{codecom}{HTML}{6E7781}

% ---------- Code / assembly listings ----------
\lstdefinestyle{c}{
  language=C,
  basicstyle=\ttfamily\small,
  keywordstyle=\color{codekw}\bfseries,
  stringstyle=\color{codestr},
  commentstyle=\color{codecom}\itshape,
  backgroundcolor=\color{codebg},
  frame=single,
  rulecolor=\color{codeframe},
  framerule=0.4pt,
  showstringspaces=false,
  breaklines=true,
  columns=fullflexible,
  keepspaces=true,
  tabsize=4,
  xleftmargin=8pt,
  xrightmargin=4pt,
  aboveskip=6pt,
  belowskip=6pt,
}
\lstdefinestyle{asm}{
  basicstyle=\ttfamily\small,
  commentstyle=\color{codecom}\itshape,
  backgroundcolor=\color{codebg},
  frame=single,
  rulecolor=\color{codeframe},
  framerule=0.4pt,
  showstringspaces=false,
  breaklines=true,
  columns=fullflexible,
  keepspaces=true,
  tabsize=4,
  xleftmargin=8pt,
  xrightmargin=4pt,
  aboveskip=6pt,
  belowskip=6pt,
}
\lstset{style=asm}

% ---------- Section styling ----------
\titleformat{\section}{\Large\bfseries\color{accent}}{\thesection.}{0.5em}{}
\titleformat{\subsection}{\large\bfseries\color{accent!85!black}}{\thesubsection}{0.5em}{}
\titlespacing*{\section}{0pt}{12pt}{6pt}
\titlespacing*{\subsection}{0pt}{8pt}{4pt}

% ---------- Page headers/footers ----------
\pagestyle{fancy}
\fancyhf{}
\fancyhead[L]{\small\CourseName\ifx\CourseCode\empty\else\ (\CourseCode)\fi}
\fancyhead[C]{\small\Semester}
\fancyhead[R]{\small\textbf{Lab \LabNumber:} \LabTitle}
\fancyfoot[C]{\small\thepage}
\renewcommand{\headrulewidth}{0.4pt}

% ---------- Title block ----------
\newcommand{\labheader}{%
  \begin{tcolorbox}[
    enhanced,
    colback=accentlight,
    colframe=accent,
    boxrule=0.5pt,
    arc=2pt,
    left=10pt, right=10pt, top=8pt, bottom=8pt
  ]
    {\Large\bfseries\color{accent} Lab \LabNumber: \LabTitle}\\[2pt]
    {\small \CourseName\ifx\CourseCode\empty\else\ (\CourseCode)\fi\quad\textbar\quad \Semester}\\[1pt]
    {\small \textbf{Duration:} \LabDuration\quad\textbar\quad \textbf{Topic:} \LabTopic}
  \end{tcolorbox}
  \vspace{2pt}
}

% ---------- Custom environments ----------
\newtcolorbox{summary}[1][Quick Reference]{
  enhanced,breakable,
  colback=accentlight!40,colframe=accent!70,
  title={\bfseries #1},
  fonttitle=\color{white},
  attach boxed title to top left={xshift=6pt,yshift=-3pt},
  boxed title style={colback=accent,boxrule=0pt,arc=2pt},
  arc=2pt,boxrule=0.4pt,
  top=12pt,
}

\newcounter{labex}[section]
\newtcolorbox{labexercise}[1]{
  enhanced,breakable,
  colback=white,colframe=accent,
  title={\bfseries In-Lab Exercise \refstepcounter{labex}\thelabex: #1},
  fonttitle=\color{white},
  attach boxed title to top left={xshift=6pt,yshift=-3pt},
  boxed title style={colback=accent,boxrule=0pt,arc=2pt},
  arc=2pt,boxrule=0.4pt,
  top=12pt,
}

\newcounter{traceex}
\newtcolorbox{trace}[1][]{
  enhanced,breakable,
  colback=white,colframe=warn,
  title={\bfseries Trace \& Predict \refstepcounter{traceex}\thetraceex\ifx\\#1\\\else: #1\fi},
  fonttitle=\color{white},
  attach boxed title to top left={xshift=6pt,yshift=-3pt},
  boxed title style={colback=warn,boxrule=0pt,arc=2pt},
  arc=2pt,boxrule=0.4pt,
  top=12pt,
}

\newcounter{bugex}
\newtcolorbox{bug}[1][]{
  enhanced,breakable,
  colback=white,colframe=warn,
  title={\bfseries Find the Bug \refstepcounter{bugex}\thebugex\ifx\\#1\\\else: #1\fi},
  fonttitle=\color{white},
  attach boxed title to top left={xshift=6pt,yshift=-3pt},
  boxed title style={colback=warn,boxrule=0pt,arc=2pt},
  arc=2pt,boxrule=0.4pt,
  top=12pt,
}

\newcounter{writeex}
\newtcolorbox{writecode}[1][]{
  enhanced,breakable,
  colback=white,colframe=warn,
  title={\bfseries Write Code on Paper \refstepcounter{writeex}\thewriteex\ifx\\#1\\\else: #1\fi},
  fonttitle=\color{white},
  attach boxed title to top left={xshift=6pt,yshift=-3pt},
  boxed title style={colback=warn,boxrule=0pt,arc=2pt},
  arc=2pt,boxrule=0.4pt,
  top=12pt,
}

% ---------- AI Collaboration box ----------
\definecolor{aipurple}{HTML}{4C1D95}
\definecolor{aipurplelight}{HTML}{EDE9FE}

\newtcolorbox{aicollab}[1][AI Collaboration]{
  enhanced,breakable,
  colback=aipurplelight,colframe=aipurple!70,
  title={\bfseries #1},
  fonttitle=\color{white},
  attach boxed title to top left={xshift=6pt,yshift=-3pt},
  boxed title style={colback=aipurple,boxrule=0pt,arc=2pt},
  arc=2pt,boxrule=0.4pt,
  top=12pt,
}

\newcounter{tryex}
\newtcolorbox{tryit}[1]{
  enhanced,breakable,
  colback=ok!6,colframe=ok!70!black,
  title={\bfseries Try It \refstepcounter{tryex}\thetryex: #1},
  fonttitle=\color{white},
  attach boxed title to top left={xshift=6pt,yshift=-3pt},
  boxed title style={colback=ok!70!black,boxrule=0pt,arc=2pt},
  arc=2pt,boxrule=0.4pt,
  top=12pt,
}

% ---------- Drawing exercise ----------
\definecolor{drawteal}{HTML}{0D7377}
\definecolor{drawteallight}{HTML}{D4F1F4}

\newcounter{drawex}
\newtcolorbox{drawex}[1]{
  enhanced,breakable,
  colback=drawteallight,colframe=drawteal!70,
  title={\bfseries Draw on Paper \refstepcounter{drawex}\thedrawex: #1},
  fonttitle=\color{white},
  attach boxed title to top left={xshift=6pt,yshift=-3pt},
  boxed title style={colback=drawteal,boxrule=0pt,arc=2pt},
  arc=2pt,boxrule=0.4pt,
  top=12pt,
}


\newcommand{\takehomebanner}{%
  \vspace{6pt}
  \begin{tcolorbox}[
    colback=warn!10,colframe=warn,boxrule=0.5pt,arc=2pt,
    left=10pt,right=10pt,top=6pt,bottom=6pt
  ]
    \textbf{Take-Home (Paper Work).}\;
    Solve the following on paper without using a computer or compiler.
    Hand-write your answers; submit at the start of the next lab.
    Goal: prepare you to dry-code during the written exam.
  \end{tcolorbox}
  \vspace{4pt}
}
 after \documentclass{article}.

\usepackage[a4paper,margin=2.2cm]{geometry}
\usepackage[T1]{fontenc}
\usepackage{lmodern}
\usepackage{microtype}
\usepackage{enumitem}
\usepackage{xcolor}
\usepackage{tcolorbox}
\tcbuselibrary{breakable,skins}
\usepackage{listings}
\usepackage{fancyhdr}
\usepackage{titlesec}
\usepackage{amsmath}
\usepackage{array}
\usepackage{booktabs}
\usepackage{hyperref}

\hypersetup{
  colorlinks=true,
  linkcolor=black,
  urlcolor=blue!60!black,
  pdfborder={0 0 0}
}

% ---------- Course identity (override per file if needed) ----------
\providecommand{\CourseName}{Microprocessors}
\providecommand{\CourseCode}{010153521}
\providecommand{\Semester}{Semester 1/2026}
\providecommand{\LabDuration}{3 hours (in-lab)}

% ---------- Colors ----------
\definecolor{accent}{HTML}{1F4E79}
\definecolor{accentlight}{HTML}{D9E2EC}
\definecolor{warn}{HTML}{B45309}
\definecolor{ok}{HTML}{166534}
\definecolor{codebg}{HTML}{F6F8FA}
\definecolor{codeframe}{HTML}{D0D7DE}
\definecolor{codekw}{HTML}{0550AE}
\definecolor{codestr}{HTML}{0A7B2C}
\definecolor{codecom}{HTML}{6E7781}

% ---------- Code / assembly listings ----------
\lstdefinestyle{c}{
  language=C,
  basicstyle=\ttfamily\small,
  keywordstyle=\color{codekw}\bfseries,
  stringstyle=\color{codestr},
  commentstyle=\color{codecom}\itshape,
  backgroundcolor=\color{codebg},
  frame=single,
  rulecolor=\color{codeframe},
  framerule=0.4pt,
  showstringspaces=false,
  breaklines=true,
  columns=fullflexible,
  keepspaces=true,
  tabsize=4,
  xleftmargin=8pt,
  xrightmargin=4pt,
  aboveskip=6pt,
  belowskip=6pt,
}
\lstdefinestyle{asm}{
  basicstyle=\ttfamily\small,
  commentstyle=\color{codecom}\itshape,
  backgroundcolor=\color{codebg},
  frame=single,
  rulecolor=\color{codeframe},
  framerule=0.4pt,
  showstringspaces=false,
  breaklines=true,
  columns=fullflexible,
  keepspaces=true,
  tabsize=4,
  xleftmargin=8pt,
  xrightmargin=4pt,
  aboveskip=6pt,
  belowskip=6pt,
}
\lstset{style=asm}

% ---------- Section styling ----------
\titleformat{\section}{\Large\bfseries\color{accent}}{\thesection.}{0.5em}{}
\titleformat{\subsection}{\large\bfseries\color{accent!85!black}}{\thesubsection}{0.5em}{}
\titlespacing*{\section}{0pt}{12pt}{6pt}
\titlespacing*{\subsection}{0pt}{8pt}{4pt}

% ---------- Page headers/footers ----------
\pagestyle{fancy}
\fancyhf{}
\fancyhead[L]{\small\CourseName\ifx\CourseCode\empty\else\ (\CourseCode)\fi}
\fancyhead[C]{\small\Semester}
\fancyhead[R]{\small\textbf{Lab \LabNumber:} \LabTitle}
\fancyfoot[C]{\small\thepage}
\renewcommand{\headrulewidth}{0.4pt}

% ---------- Title block ----------
\newcommand{\labheader}{%
  \begin{tcolorbox}[
    enhanced,
    colback=accentlight,
    colframe=accent,
    boxrule=0.5pt,
    arc=2pt,
    left=10pt, right=10pt, top=8pt, bottom=8pt
  ]
    {\Large\bfseries\color{accent} Lab \LabNumber: \LabTitle}\\[2pt]
    {\small \CourseName\ifx\CourseCode\empty\else\ (\CourseCode)\fi\quad\textbar\quad \Semester}\\[1pt]
    {\small \textbf{Duration:} \LabDuration\quad\textbar\quad \textbf{Topic:} \LabTopic}
  \end{tcolorbox}
  \vspace{2pt}
}

% ---------- Custom environments ----------
\newtcolorbox{summary}[1][Quick Reference]{
  enhanced,breakable,
  colback=accentlight!40,colframe=accent!70,
  title={\bfseries #1},
  fonttitle=\color{white},
  attach boxed title to top left={xshift=6pt,yshift=-3pt},
  boxed title style={colback=accent,boxrule=0pt,arc=2pt},
  arc=2pt,boxrule=0.4pt,
  top=12pt,
}

\newcounter{labex}[section]
\newtcolorbox{labexercise}[1]{
  enhanced,breakable,
  colback=white,colframe=accent,
  title={\bfseries In-Lab Exercise \refstepcounter{labex}\thelabex: #1},
  fonttitle=\color{white},
  attach boxed title to top left={xshift=6pt,yshift=-3pt},
  boxed title style={colback=accent,boxrule=0pt,arc=2pt},
  arc=2pt,boxrule=0.4pt,
  top=12pt,
}

\newcounter{traceex}
\newtcolorbox{trace}[1][]{
  enhanced,breakable,
  colback=white,colframe=warn,
  title={\bfseries Trace \& Predict \refstepcounter{traceex}\thetraceex\ifx\\#1\\\else: #1\fi},
  fonttitle=\color{white},
  attach boxed title to top left={xshift=6pt,yshift=-3pt},
  boxed title style={colback=warn,boxrule=0pt,arc=2pt},
  arc=2pt,boxrule=0.4pt,
  top=12pt,
}

\newcounter{bugex}
\newtcolorbox{bug}[1][]{
  enhanced,breakable,
  colback=white,colframe=warn,
  title={\bfseries Find the Bug \refstepcounter{bugex}\thebugex\ifx\\#1\\\else: #1\fi},
  fonttitle=\color{white},
  attach boxed title to top left={xshift=6pt,yshift=-3pt},
  boxed title style={colback=warn,boxrule=0pt,arc=2pt},
  arc=2pt,boxrule=0.4pt,
  top=12pt,
}

\newcounter{writeex}
\newtcolorbox{writecode}[1][]{
  enhanced,breakable,
  colback=white,colframe=warn,
  title={\bfseries Write Code on Paper \refstepcounter{writeex}\thewriteex\ifx\\#1\\\else: #1\fi},
  fonttitle=\color{white},
  attach boxed title to top left={xshift=6pt,yshift=-3pt},
  boxed title style={colback=warn,boxrule=0pt,arc=2pt},
  arc=2pt,boxrule=0.4pt,
  top=12pt,
}

% ---------- AI Collaboration box ----------
\definecolor{aipurple}{HTML}{4C1D95}
\definecolor{aipurplelight}{HTML}{EDE9FE}

\newtcolorbox{aicollab}[1][AI Collaboration]{
  enhanced,breakable,
  colback=aipurplelight,colframe=aipurple!70,
  title={\bfseries #1},
  fonttitle=\color{white},
  attach boxed title to top left={xshift=6pt,yshift=-3pt},
  boxed title style={colback=aipurple,boxrule=0pt,arc=2pt},
  arc=2pt,boxrule=0.4pt,
  top=12pt,
}

\newcounter{tryex}
\newtcolorbox{tryit}[1]{
  enhanced,breakable,
  colback=ok!6,colframe=ok!70!black,
  title={\bfseries Try It \refstepcounter{tryex}\thetryex: #1},
  fonttitle=\color{white},
  attach boxed title to top left={xshift=6pt,yshift=-3pt},
  boxed title style={colback=ok!70!black,boxrule=0pt,arc=2pt},
  arc=2pt,boxrule=0.4pt,
  top=12pt,
}

% ---------- Drawing exercise ----------
\definecolor{drawteal}{HTML}{0D7377}
\definecolor{drawteallight}{HTML}{D4F1F4}

\newcounter{drawex}
\newtcolorbox{drawex}[1]{
  enhanced,breakable,
  colback=drawteallight,colframe=drawteal!70,
  title={\bfseries Draw on Paper \refstepcounter{drawex}\thedrawex: #1},
  fonttitle=\color{white},
  attach boxed title to top left={xshift=6pt,yshift=-3pt},
  boxed title style={colback=drawteal,boxrule=0pt,arc=2pt},
  arc=2pt,boxrule=0.4pt,
  top=12pt,
}


\newcommand{\takehomebanner}{%
  \vspace{6pt}
  \begin{tcolorbox}[
    colback=warn!10,colframe=warn,boxrule=0.5pt,arc=2pt,
    left=10pt,right=10pt,top=6pt,bottom=6pt
  ]
    \textbf{Take-Home (Paper Work).}\;
    Solve the following on paper without using a computer or compiler.
    Hand-write your answers; submit at the start of the next lab.
    Goal: prepare you to dry-code during the written exam.
  \end{tcolorbox}
  \vspace{4pt}
}
 after \documentclass{article}.

\usepackage[a4paper,margin=2.2cm]{geometry}
\usepackage[T1]{fontenc}
\usepackage{lmodern}
\usepackage{microtype}
\usepackage{enumitem}
\usepackage{xcolor}
\usepackage{tcolorbox}
\tcbuselibrary{breakable,skins}
\usepackage{listings}
\usepackage{fancyhdr}
\usepackage{titlesec}
\usepackage{amsmath}
\usepackage{array}
\usepackage{booktabs}
\usepackage{hyperref}

\hypersetup{
  colorlinks=true,
  linkcolor=black,
  urlcolor=blue!60!black,
  pdfborder={0 0 0}
}

% ---------- Course identity (override per file if needed) ----------
\providecommand{\CourseName}{Microprocessors}
\providecommand{\CourseCode}{010153521}
\providecommand{\Semester}{Semester 1/2026}
\providecommand{\LabDuration}{3 hours (in-lab)}

% ---------- Colors ----------
\definecolor{accent}{HTML}{1F4E79}
\definecolor{accentlight}{HTML}{D9E2EC}
\definecolor{warn}{HTML}{B45309}
\definecolor{ok}{HTML}{166534}
\definecolor{codebg}{HTML}{F6F8FA}
\definecolor{codeframe}{HTML}{D0D7DE}
\definecolor{codekw}{HTML}{0550AE}
\definecolor{codestr}{HTML}{0A7B2C}
\definecolor{codecom}{HTML}{6E7781}

% ---------- Code / assembly listings ----------
\lstdefinestyle{c}{
  language=C,
  basicstyle=\ttfamily\small,
  keywordstyle=\color{codekw}\bfseries,
  stringstyle=\color{codestr},
  commentstyle=\color{codecom}\itshape,
  backgroundcolor=\color{codebg},
  frame=single,
  rulecolor=\color{codeframe},
  framerule=0.4pt,
  showstringspaces=false,
  breaklines=true,
  columns=fullflexible,
  keepspaces=true,
  tabsize=4,
  xleftmargin=8pt,
  xrightmargin=4pt,
  aboveskip=6pt,
  belowskip=6pt,
}
\lstdefinestyle{asm}{
  basicstyle=\ttfamily\small,
  commentstyle=\color{codecom}\itshape,
  backgroundcolor=\color{codebg},
  frame=single,
  rulecolor=\color{codeframe},
  framerule=0.4pt,
  showstringspaces=false,
  breaklines=true,
  columns=fullflexible,
  keepspaces=true,
  tabsize=4,
  xleftmargin=8pt,
  xrightmargin=4pt,
  aboveskip=6pt,
  belowskip=6pt,
}
\lstset{style=asm}

% ---------- Section styling ----------
\titleformat{\section}{\Large\bfseries\color{accent}}{\thesection.}{0.5em}{}
\titleformat{\subsection}{\large\bfseries\color{accent!85!black}}{\thesubsection}{0.5em}{}
\titlespacing*{\section}{0pt}{12pt}{6pt}
\titlespacing*{\subsection}{0pt}{8pt}{4pt}

% ---------- Page headers/footers ----------
\pagestyle{fancy}
\fancyhf{}
\fancyhead[L]{\small\CourseName\ifx\CourseCode\empty\else\ (\CourseCode)\fi}
\fancyhead[C]{\small\Semester}
\fancyhead[R]{\small\textbf{Lab \LabNumber:} \LabTitle}
\fancyfoot[C]{\small\thepage}
\renewcommand{\headrulewidth}{0.4pt}

% ---------- Title block ----------
\newcommand{\labheader}{%
  \begin{tcolorbox}[
    enhanced,
    colback=accentlight,
    colframe=accent,
    boxrule=0.5pt,
    arc=2pt,
    left=10pt, right=10pt, top=8pt, bottom=8pt
  ]
    {\Large\bfseries\color{accent} Lab \LabNumber: \LabTitle}\\[2pt]
    {\small \CourseName\ifx\CourseCode\empty\else\ (\CourseCode)\fi\quad\textbar\quad \Semester}\\[1pt]
    {\small \textbf{Duration:} \LabDuration\quad\textbar\quad \textbf{Topic:} \LabTopic}
  \end{tcolorbox}
  \vspace{2pt}
}

% ---------- Custom environments ----------
\newtcolorbox{summary}[1][Quick Reference]{
  enhanced,breakable,
  colback=accentlight!40,colframe=accent!70,
  title={\bfseries #1},
  fonttitle=\color{white},
  attach boxed title to top left={xshift=6pt,yshift=-3pt},
  boxed title style={colback=accent,boxrule=0pt,arc=2pt},
  arc=2pt,boxrule=0.4pt,
  top=12pt,
}

\newcounter{labex}[section]
\newtcolorbox{labexercise}[1]{
  enhanced,breakable,
  colback=white,colframe=accent,
  title={\bfseries In-Lab Exercise \refstepcounter{labex}\thelabex: #1},
  fonttitle=\color{white},
  attach boxed title to top left={xshift=6pt,yshift=-3pt},
  boxed title style={colback=accent,boxrule=0pt,arc=2pt},
  arc=2pt,boxrule=0.4pt,
  top=12pt,
}

\newcounter{traceex}
\newtcolorbox{trace}[1][]{
  enhanced,breakable,
  colback=white,colframe=warn,
  title={\bfseries Trace \& Predict \refstepcounter{traceex}\thetraceex\ifx\\#1\\\else: #1\fi},
  fonttitle=\color{white},
  attach boxed title to top left={xshift=6pt,yshift=-3pt},
  boxed title style={colback=warn,boxrule=0pt,arc=2pt},
  arc=2pt,boxrule=0.4pt,
  top=12pt,
}

\newcounter{bugex}
\newtcolorbox{bug}[1][]{
  enhanced,breakable,
  colback=white,colframe=warn,
  title={\bfseries Find the Bug \refstepcounter{bugex}\thebugex\ifx\\#1\\\else: #1\fi},
  fonttitle=\color{white},
  attach boxed title to top left={xshift=6pt,yshift=-3pt},
  boxed title style={colback=warn,boxrule=0pt,arc=2pt},
  arc=2pt,boxrule=0.4pt,
  top=12pt,
}

\newcounter{writeex}
\newtcolorbox{writecode}[1][]{
  enhanced,breakable,
  colback=white,colframe=warn,
  title={\bfseries Write Code on Paper \refstepcounter{writeex}\thewriteex\ifx\\#1\\\else: #1\fi},
  fonttitle=\color{white},
  attach boxed title to top left={xshift=6pt,yshift=-3pt},
  boxed title style={colback=warn,boxrule=0pt,arc=2pt},
  arc=2pt,boxrule=0.4pt,
  top=12pt,
}

% ---------- AI Collaboration box ----------
\definecolor{aipurple}{HTML}{4C1D95}
\definecolor{aipurplelight}{HTML}{EDE9FE}

\newtcolorbox{aicollab}[1][AI Collaboration]{
  enhanced,breakable,
  colback=aipurplelight,colframe=aipurple!70,
  title={\bfseries #1},
  fonttitle=\color{white},
  attach boxed title to top left={xshift=6pt,yshift=-3pt},
  boxed title style={colback=aipurple,boxrule=0pt,arc=2pt},
  arc=2pt,boxrule=0.4pt,
  top=12pt,
}

\newcounter{tryex}
\newtcolorbox{tryit}[1]{
  enhanced,breakable,
  colback=ok!6,colframe=ok!70!black,
  title={\bfseries Try It \refstepcounter{tryex}\thetryex: #1},
  fonttitle=\color{white},
  attach boxed title to top left={xshift=6pt,yshift=-3pt},
  boxed title style={colback=ok!70!black,boxrule=0pt,arc=2pt},
  arc=2pt,boxrule=0.4pt,
  top=12pt,
}

% ---------- Drawing exercise ----------
\definecolor{drawteal}{HTML}{0D7377}
\definecolor{drawteallight}{HTML}{D4F1F4}

\newcounter{drawex}
\newtcolorbox{drawex}[1]{
  enhanced,breakable,
  colback=drawteallight,colframe=drawteal!70,
  title={\bfseries Draw on Paper \refstepcounter{drawex}\thedrawex: #1},
  fonttitle=\color{white},
  attach boxed title to top left={xshift=6pt,yshift=-3pt},
  boxed title style={colback=drawteal,boxrule=0pt,arc=2pt},
  arc=2pt,boxrule=0.4pt,
  top=12pt,
}


\newcommand{\takehomebanner}{%
  \vspace{6pt}
  \begin{tcolorbox}[
    colback=warn!10,colframe=warn,boxrule=0.5pt,arc=2pt,
    left=10pt,right=10pt,top=6pt,bottom=6pt
  ]
    \textbf{Take-Home (Paper Work).}\;
    Solve the following on paper without using a computer or compiler.
    Hand-write your answers; submit at the start of the next lab.
    Goal: prepare you to dry-code during the written exam.
  \end{tcolorbox}
  \vspace{4pt}
}


\begin{document}
\labheader

\section*{Objectives}
\begin{itemize}[leftmargin=*,nosep]
  \item Configure a real I/O pin as an output with \texttt{DDRB} and drive it with
        \texttt{OUT} -- direct I/O addressing on physical hardware, not a simulator
        (Lecture 2, \S2.2-2.3).
  \item Write a correct, buildable AVR assembly program in Microchip Studio: correct
        program structure, \texttt{.INCLUDE}, \texttt{.ORG}, and an initialized stack
        pointer before any \texttt{CALL}/\texttt{RET} (Lecture 2, \S2.5-2.6; Lecture 3,
        \S3.2).
  \item Build single- and nested-counter delay loops with \texttt{DEC}/\texttt{BRNE},
        wrapped in a \texttt{CALL}/\texttt{RET} subroutine (Lecture 3, \S3.1-3.2).
  \item Calculate, by hand, the number of loop iterations needed for a target time
        delay at the Arduino Uno's actual crystal frequency (16 MHz), using the
        cycle-counting method of Lecture 3, \S3.3.
  \item Assemble the program to an Intel HEX file and upload it to an Arduino Uno R3
        using \texttt{avrdude}'s Arduino (bootloader) programmer -- entirely from the
        command line, without the Arduino IDE.
  \item Verify the real output waveform on an oscilloscope -- period, frequency, duty
        cycle -- against your hand calculation, and account for the difference.
\end{itemize}

\begin{summary}[Quick Reference: I/O on This Board]
The ATmega328P's data memory, SFR names, and \texttt{IN}/\texttt{OUT} mechanics are
exactly as in Lecture 2 -- only the chip and the pin names change.
\begin{center}
\begin{tabular}{@{}p{3.7cm}p{1.7cm}p{8.9cm}@{}}
\toprule
\textbf{Port} & \textbf{Dir. reg.} & \textbf{Arduino Uno pins} \\
\midrule
PORTB (\texttt{PB0-PB5}) & \texttt{DDRB} & D8-D13 (\texttt{PB5} = D13, on-board LED) \\
PORTC (\texttt{PC0-PC5}) & \texttt{DDRC} & A0-A5 \\
PORTD (\texttt{PD0-PD7}) & \texttt{DDRD} & D0-D7 (\texttt{PD0}/\texttt{PD1} = USB
\textbf{RXD}/\textbf{TXD} -- do not drive) \\
\bottomrule
\end{tabular}
\end{center}
\texttt{DDRx} bit = 1 makes that pin an \emph{output}; bit = 0 leaves it an
\emph{input} (0 is the reset default, i.e. \emph{every} pin starts as an input).
\texttt{OUT PORTx,Rr} only \emph{drives} a pin once its \texttt{DDRx} bit is set --
otherwise it just toggles that pin's internal pull-up (see the ``Find the Bug''
exercise). Register names such as \texttt{PB5} and \texttt{DDRB} come from
\texttt{m328Pdef.inc}, exactly as \texttt{M32DEF.INC} supplied \texttt{PORTB} in
Lecture 2's Program 2-1.
\end{summary}

\begin{summary}[Quick Reference: Delay Math at 16 MHz]
The Arduino Uno's crystal is \textbf{16 MHz}, so
$T_{\text{cyc}} = 1/16\text{MHz} = 62.5$~ns -- half the cycle time used in Lecture 3's
worked examples (mostly 1 MHz or 8 MHz). The delay formula is unchanged:
\[
\text{delay} = \big[\,1 + (C \times N) - 1 + 4\,\big] \times T_{\text{cyc}}
\]
where $N$ is the loop count (counter register), $C$ is the cycles used per pass
\emph{including} the branch when it is taken, the leading 1 accounts for loading the
counter, the $-1$ corrects for the final (non-taken) branch, and the 4 accounts for
\texttt{RET}.

\textbf{Worked demo (a different target from any lab exercise) -- aim for
100~\textmu s:} target cycles $= 100{,}000\text{ns} / 62.5\text{ns} = 1600$. Choose
$N=200$ and pad the loop body with 5 \texttt{NOP}s, so
$C = 5(\texttt{NOP}) + 1(\texttt{DEC}) + 2(\texttt{BRNE, taken}) = 8$ cycles/pass.
\[
\big[\,1 + (8\times200) - 1 + 4\,\big]\times62.5\text{ns} = 1604 \times 62.5\text{ns}
= 100.25\ \text{\textmu s} \qquad (\text{error} = +0.25\%)
\]
Perfect round numbers are rare -- getting within a percent or two of the target and
then \emph{measuring} the real error on the scope is the point of this lab, not
hitting the target exactly (see ``Verifying Delays in Practice,'' Lecture 3).
\end{summary}

\begin{summary}[Quick Reference: Your Individual Target]
Exercises 2 and 3 ask for \emph{your own} target, not a shared one -- so a copied
\texttt{.asm} won't hit the right numbers. Derive it from your student ID with this
hash:
\[
D = \text{sum of every digit in your student ID}
\]
\begin{center}
\begin{tabular}{@{}p{5.7cm}p{8.9cm}@{}}
\toprule
\textbf{Exercise} & \textbf{Your target} \\
\midrule
2 (single-loop delay) & $T_2 = 0.50\text{ ms} + (D \bmod 20) \times 0.05\text{ ms}$ \\
3 (nested-loop square wave) & $f_3 = 500\text{ Hz} + (D \bmod 30) \times 50\text{ Hz}$ \\
\bottomrule
\end{tabular}
\end{center}
\textbf{Worked example} (a made-up 13-digit ID, same format as yours -- do not reuse
these numbers): ID \texttt{1029384756123} $\to$
$D = 1{+}0{+}2{+}9{+}3{+}8{+}4{+}7{+}5{+}6{+}1{+}2{+}3 = 51$. Then
$T_2 = 0.50 + (51 \bmod 20)\times0.05 = 0.50 + 11\times0.05 = 1.05$~ms, and
$f_3 = 500 + (51 \bmod 30)\times50 = 500 + 21\times50 = 1550$~Hz.

Compute your own $D$, $T_2$, and $f_3$ now and record them at the top of your lab
notebook -- your instructor will spot-check that your measured results match
\emph{your} numbers. Working in pairs? Use whichever partner's ID is listed first on
your roster.
\end{summary}

\section*{Toolchain: From Source Code to Oscilloscope Trace}

\begin{enumerate}[leftmargin=*,nosep]
  \item Write the \texttt{.asm} source in \textbf{Microchip Studio} (Assembler
        project, device \texttt{ATmega328P}).
  \item \textbf{Build} (F7). Studio's assembler produces the usual family of
        output files (\S2.7) in the project's \texttt{Debug} folder, including a
        ready-to-burn \texttt{.hex} file and a \texttt{.lst} listing.
  \item \textbf{(Recommended) Simulate} the \texttt{.hex} in an online simulator
        such as \textbf{Wokwi} to catch logic mistakes -- \emph{not} timing --
        before touching the real board.
  \item Upload the \texttt{.hex} file to the board with \textbf{\texttt{avrdude}},
        run from a terminal -- \emph{not} Studio's built-in Device Programming
        dialog, which targets ISP/debugWIRE hardware programmers, not the Uno's
        onboard USB-serial bootloader.
  \item Probe the target pin with an \textbf{oscilloscope} and compare the measured
        waveform against your hand calculation.
\end{enumerate}

\subsection*{Setting Up the Microchip Studio Project}
\begin{enumerate}[leftmargin=*,nosep]
  \item \textbf{File > New > Project... > Assembler > AVR Assembler.} Name the
        project (e.g. \texttt{Lab3\_Delay}).
  \item On the device-selection dialog, choose \textbf{ATmega328P} (\emph{not}
        ATmega328 -- the ``P'' variant is what ships on the Uno R3).
  \item Replace the generated \texttt{main.asm} with the skeleton in
        Listing~\ref{lst:skeleton} below.
\end{enumerate}

\subsection*{The Skeleton Program}

\begin{lstlisting}[caption={Skeleton: blink-check on PB5 (Arduino D13). The delay is
a placeholder -- not yet timed.},label={lst:skeleton}]
;Lab 3 skeleton -- Arduino Uno R3 (ATmega328P), 16 MHz crystal

.INCLUDE "m328Pdef.inc"

.ORG 0x0000
        LDI   R16, HIGH(RAMEND)   ;initialize the stack pointer --
        OUT   SPH, R16            ;required before any CALL/RCALL/RET
        LDI   R16, LOW(RAMEND)
        OUT   SPL, R16

        LDI   R16, (1<<PB5)       ;PB5 = output (Arduino D13), all other
        OUT   DDRB, R16           ;PORTB pins left as inputs

MAIN:   LDI   R16, (1<<PB5)       ;drive PB5 high
        OUT   PORTB, R16
        RCALL DELAY
        LDI   R16, 0x00           ;drive PB5 low
        OUT   PORTB, R16
        RCALL DELAY
        RJMP  MAIN

;----- delay subroutine (placeholder -- retime this in Exercise 2) -----
DELAY:  LDI   R20, 0xFF
AGAIN:  DEC   R20
        BRNE  AGAIN
        RET
\end{lstlisting}

* We do \emph{not} reserve interrupt vector slots: no interrupt is ever enabled in
this program, so \texttt{.ORG 0x0000} places our first instruction directly at the
reset vector -- exactly like Program 2-1's \texttt{.ORG 00}.

\subsection*{Simulate First (Recommended): Wokwi}

Before burning a real chip, it is worth checking your program's \emph{logic} in a
free online simulator such as \textbf{Wokwi} (wokwi.com):
\begin{enumerate}[leftmargin=*,nosep]
  \item Start a new Wokwi \textbf{Arduino Uno} project.
  \item Upload the \texttt{.hex} file Microchip Studio just built as the project's
        firmware -- Wokwi's Arduino Uno will run any valid AVR \texttt{.hex}, whether
        or not Wokwi itself compiled it. No Arduino sketch is needed.
  \item Wire a virtual LED (or Wokwi's logic analyzer) to the pin under test and
        press Run.
\end{enumerate}
This catches the mistakes that have nothing to do with timing -- a missing
\texttt{DDRB} write, a typo'd register, a loop that never exits -- instantly and for
free, with no board plugged in.

\textbf{What Wokwi cannot tell you:} its simulated clock is \emph{not} cycle-accurate
to a real 16 MHz crystal, so it cannot confirm whether your delay actually hits your
target $T_2$ or $f_3$. Use it only to confirm the waveform's \emph{shape} -- does the
pin toggle at all, does the counter pattern look right -- never to answer this lab's
timing questions. Only the physical oscilloscope measurement counts for that.

\subsection*{Uploading with AVRdude}

\texttt{avrdude} talks to the Uno's factory-installed bootloader over the same USB
cable used for programming from the Arduino IDE -- the \texttt{arduino} programmer ID
selects the STK500v1 protocol that bootloader speaks, at 115200 baud.

\begin{lstlisting}[language=bash]
avrdude -c arduino -p atmega328p -P <port> -b 115200 -U flash:w:Lab3_Delay.hex:i
\end{lstlisting}

\begin{center}
\begin{tabular}{@{}p{4.2cm}p{8.2cm}@{}}
\toprule
\textbf{Flag} & \textbf{Meaning} \\
\midrule
\texttt{-c arduino} & Programmer type: Arduino/Optiboot bootloader (STK500v1). \\
\texttt{-p atmega328p} & Target part. \\
\texttt{-P <port>} & Serial port: \texttt{COM5} (Windows), \texttt{/dev/cu.usbmodemXXXX}
(macOS), \texttt{/dev/ttyACM0} (Linux). Some Uno clones use a CH340 USB-serial chip
and need a separate driver before the port appears. \\
\texttt{-b 115200} & Bootloader baud rate (Uno-specific; older boards use 19200). \\
\texttt{-U flash:w:} \texttt{file.hex:i} & Write \texttt{file.hex} (Intel HEX format)
to flash. \\
\bottomrule
\end{tabular}
\end{center}

\textbf{Before uploading:} close the Arduino IDE's Serial Monitor and any terminal
program holding the port open -- \texttt{avrdude} needs exclusive access, and a busy
port is the most common upload failure (Appendix A).

\textbf{Never} pass \texttt{-U lfuse:w:...}, \texttt{hfuse}, or \texttt{efuse} on this
board unless explicitly instructed to. The Uno's fuses are what make the bootloader
run at power-up; writing the wrong fuse value can make the board unreachable by
\texttt{avrdude} (a full recovery then needs a hardware ISP programmer).

\subsection*{Connecting the Oscilloscope}
\begin{itemize}[leftmargin=*,nosep]
  \item Probe tip on the pin under test (e.g. the D13 header pin); probe ground clip
        on any Uno \textbf{GND} pin. A shared ground between board and scope is
        required for a clean trace.
  \item Start with the vertical scale at 1-2 V/div (logic swings 0-5 V) and the
        timebase set so that 2-3 full periods fit on screen; then use the scope's
        auto-measure or cursors to read period, frequency, and duty cycle.
  \item If the trace looks noisy or doesn't trigger, check that \texttt{DDRB} was
        actually set for that pin (see the ``Find the Bug'' exercise) before
        suspecting the scope.
\end{itemize}

\subsection*{Quick Experiments (Try It)}

\begin{tryit}{Hello, pin 13}
Build the skeleton exactly as given. If you like, load the \texttt{.hex} into Wokwi
first and confirm the virtual pin toggles -- a quick, board-free sanity check on the
\emph{logic}. Then upload the real \texttt{.hex} with \texttt{avrdude} and probe D13.
The placeholder delay is deliberately \emph{not} timed to anything meaningful --
you are not trying to see it blink. Confirm only that the scope shows \emph{a}
periodic square wave: this is your checkpoint that editor $\to$ assembler $\to$
\texttt{avrdude} $\to$ real hardware all works, before any precision math begins.
\end{tryit}

\begin{tryit}{Read the .lst file}
Open the project's \texttt{.lst} listing file (Studio writes it next to the
\texttt{.hex}) and find the machine code generated for the \texttt{RCALL DELAY} and
\texttt{BRNE AGAIN} lines.
\begin{enumerate}[nosep,leftmargin=*]
  \item Is \texttt{RCALL} 2 bytes or 4, in this listing? Does that match \S3.1's
        description of \texttt{RCALL} versus \texttt{CALL}?
  \item Using the addresses shown in the listing (not the source line numbers),
        recompute \texttt{BRNE AGAIN}'s relative displacement \texttt{k} by hand, the
        same way Example 3-7 did. Does it match the machine code you see?
\end{enumerate}
\end{tryit}

\section*{In-Lab Exercises (target: 3 hours)}

\begin{labexercise}{Toolchain checkpoint \hfill {\small \itshape (25 min)}}
Repeat Try It 1 as a graded checkpoint: project set up for \texttt{ATmega328P},
skeleton typed in exactly, built, and uploaded with \texttt{avrdude}. Record the
exact \texttt{avrdude} command you used (including your port name) and a screenshot
or sketch of the oscilloscope trace on D13. If \texttt{avrdude} fails to connect,
work through Appendix A before asking for help.
\end{labexercise}

\begin{labexercise}{Your precisely timed interval, $T_2$ \hfill {\small \itshape (35 min)}}
Using the ``Your Individual Target'' box above, compute your own $D$ and $T_2$ if you
haven't already. Replace the placeholder \texttt{DELAY} subroutine with a
\emph{single}-counter loop (counter register + \texttt{NOP} padding +
\texttt{DEC}/\texttt{BRNE}), following the method of the 100~\textmu s worked demo
above, tuned for a delay as close as possible to \textbf{your $T_2$} at 16 MHz. Show
your full cycle accounting ($N$, number of \texttt{NOP}s, $C$, total cycles, total
time, and the \% error against $T_2$) in your lab notebook.
\begin{enumerate}[nosep,leftmargin=*]
  \item Wire it into the skeleton's \texttt{MAIN} toggle loop -- with two calls to
        \texttt{DELAY} per period, PB5 should now produce a square wave with a
        period of \textbf{$2\times T_2$} (frequency $1/(2T_2)$).
  \item Build, upload, and measure the actual period and frequency on the
        oscilloscope. Compute the \% error between your \emph{hand-calculated} period
        and the \emph{measured} one, and compare it to the \% error you calculated
        against $T_2$ above -- are they close? If not, what did your hand
        calculation leave out?
\end{enumerate}
\end{labexercise}

\begin{labexercise}{Your square wave with a nested loop, $f_3$ \hfill {\small \itshape (55 min)}}
A single 8-bit counter tops out at 255 passes -- not enough range for finer control
without absurd \texttt{NOP} padding. Using your $f_3$ from the box above, design a
\emph{nested} two-register loop (as in Lecture 3's Example 3-18) for a half-period as
close as possible to \textbf{$1/(2f_3)$}, giving an overall square wave at
\textbf{your $f_3$}, roughly 50\% duty-cycle, on PB5.
\begin{enumerate}[nosep,leftmargin=*]
  \item Calculate your two register values ($N_1$, $N_2$) and the resulting delay by
        hand before touching the keyboard.
  \item Build, upload, and measure the frequency on the oscilloscope.
  \item Tune your register constants and rebuild until the \emph{measured} frequency
        is within $\pm 2\%$ of $f_3$ -- a realistic engineering tolerance, not exact
        equality. Keep a small table in your notebook: attempt \#, $N_1$, $N_2$,
        calculated period, measured period, \% error.
\end{enumerate}
\end{labexercise}

\begin{labexercise}{Binary ripple counter on PORTD \hfill {\small \itshape (35 min)}}
Combine \S2.2-2.3's \texttt{OUT} addressing with \S3.1-3.3's counting and delay
techniques: maintain an 8-bit counter register, \texttt{INC} it once per step, shift
it left \emph{twice} with \texttt{LSL} before each \texttt{OUT PORTD,Rd}, and hold it
for a fixed per-step delay (reuse or rescale your \texttt{DELAY} subroutine) before
incrementing again.
\begin{itemize}[nosep,leftmargin=*]
  \item The two \texttt{LSL}s matter: they guarantee bits \texttt{PD0}/\texttt{PD1}
        stay 0 for every count, so the USB-serial lines are never driven while the
        board stays connected -- only \texttt{PD2} and above carry the counting
        pattern.
  \item Probe \textbf{PD2} (the fastest-toggling driven bit). Its period should equal
        exactly \textbf{twice} your programmed per-step delay -- explain why in terms
        of how an incrementing binary count flips its LSB.
  \item Probe \textbf{PD3}. Its period should be exactly \textbf{twice} \texttt{PD2}'s
        -- verifying, directly on the scope, the binary-counting relationship between
        adjacent bits.
\end{itemize}
\end{labexercise}

\begin{aicollab}
\textbf{Try these prompts with an AI tool:}
\begin{itemize}[leftmargin=*,nosep]
  \item ``I need a nested AVR delay loop as close as possible to my assigned
        half-period, $1/(2f_3)$, at 16 MHz. Show the cycle-counting, not just a
        guessed answer.''
  \item ``What does \texttt{avrdude: stk500\_getsync(): not in sync} mean, and what
        are the most common causes on an Arduino Uno?''
  \item ``What is the difference between setting a \texttt{DDRx} bit and setting the
        corresponding \texttt{PORTx} bit on an AVR pin configured as an input?''
\end{itemize}

\medskip\textbf{Critical check -- verify, don't just accept:}
If an AI tool gives you register values for a target delay, do \emph{not} trust the
number -- redo its cycle count by hand, instruction by instruction, the same way you
did for the 100~\textmu s demo. Then check it a second way: does the
\emph{oscilloscope} agree, within a percent or two, with either calculation?

\medskip\textbf{Know the limit:}
AI tools are good at producing \emph{a} delay loop that is plausible-looking but will
happily miscount the cycle cost of the final (non-taken) branch, or forget the
\texttt{RET} overhead entirely -- exactly the kind of off-by-a-few-cycles error that
is invisible in the source but visible on a scope trace. The oscilloscope is your
ground truth in this lab, not the AI's arithmetic.
\end{aicollab}

\takehomebanner

\begin{trace}[Nested-loop delay at 8 MHz]
For an ATmega32 (a different chip from this lab, running at 8 MHz, so
$T_{\text{cyc}} = 125$~ns), a nested-loop \texttt{DELAY} subroutine uses
$N_1 = 20$ (outer) and $N_2 = 150$ (inner), with a 3-cycle inner loop body
(\texttt{NOP} + \texttt{DEC} + \texttt{BRNE}). Calculate the total delay by hand,
following the same two-level accounting Example 3-18 used, and give your answer in
microseconds.
\end{trace}

\begin{bug}[DDRB never set]
A student wrote the skeleton's \texttt{MAIN} loop and \texttt{DELAY} subroutine
correctly, but accidentally deleted the two lines that set \texttt{DDRB} before
\texttt{MAIN}. The program assembles without errors, and \texttt{avrdude} reports a
successful upload with no verification failure. Predict, qualitatively, what the LED
and the oscilloscope trace on D13 will show, and explain \emph{why} -- in terms of
what a \texttt{DDRx} bit versus a \texttt{PORTx} bit actually controls on a pin that
is still configured as an input -- even though every instruction executed exactly as
written.
\end{bug}

\begin{writecode}[A delay subroutine for a different chip]
On paper, without a computer, simulator, or AI tool, write a complete
\texttt{DELAY} subroutine (the two-part structure: set counter(s), then loop) that
produces as close as possible to \textbf{20 ms} on a chip running at a
\textbf{4 MHz} crystal. Show every step of your cycle accounting -- the number of
registers you need, their values, the cycles per pass, and the final delay in
microseconds compared to the 20 ms target.
\end{writecode}

\begin{drawex}{Predict, then verify, your $f_3$ trace}
Before building Exercise 3's program, sketch on paper the oscilloscope trace you
expect: axes labeled (voltage on $y$, time on $x$), the two logic levels marked
(0 V / 5 V), and the period, half-period, and duty cycle labeled with numbers.
After you build, upload, and actually measure it, annotate your sketch with the real
measured values next to your predicted ones. Are they close? If your predicted duty
cycle wasn't exactly 50\%, can you identify which part of the code caused the
asymmetry?
\end{drawex}

\section*{Reflection}

In one short paragraph, compare this lab's timing model to Lab 1/2's EC-1 machine.
On EC-1, every state transition took \emph{exactly} one clock edge -- timing was
deterministic the moment you designed the FSM, with no measurement needed to know how
long anything took. Here, the ATmega328P's crystal sets a real physical time base,
but your hand-calculated delay is still only an \emph{approximation} (it ignores
setup overhead, and your crystal has some tolerance) -- which is why this lab ends
with an oscilloscope, not just a calculation. What would you need to add to this
program to make its timing as deterministic as EC-1's, without touching the crystal
itself? (You are not expected to implement it -- Lecture 3's closing note names the
AVR feature that does exactly this, covered fully in a later chapter.)

\newpage
\section*{Appendix A: AVRdude Troubleshooting}

\begin{center}
\begin{tabular}{@{}p{5.4cm}p{6.8cm}@{}}
\toprule
\textbf{Message / symptom} & \textbf{Likely cause and fix} \\
\midrule
\texttt{stk500\_getsync(): attempt N of 10: not in sync} &
Wrong port, or the bootloader isn't listening. Double-check \texttt{-P}; press the
Uno's reset button just before running \texttt{avrdude} if auto-reset (via DTR)
isn't working. \\
\addlinespace
Port open/permission error (e.g. \texttt{Resource busy}, \texttt{Access is denied}) &
Another program (Arduino IDE Serial Monitor, a terminal, Microchip Studio's own
terminal) is holding the port open. Close it and retry. \\
\addlinespace
Port does not appear in the OS's device list at all &
Driver missing for a CH340/CH341 clone USB-serial chip; install the vendor driver.
On macOS/Linux, list ports with \texttt{ls /dev/cu.*} or \texttt{ls /dev/tty*} before
and after plugging in the board to spot it. \\
\addlinespace
\texttt{avrdude: verification error, first mismatch at byte 0x0000} &
Flash didn't actually update -- often a stale \texttt{.hex} from a failed rebuild.
Rebuild in Studio, confirm the \texttt{.hex} timestamp, and re-upload. \\
\addlinespace
Upload succeeds, but the board does nothing new &
Confirm you're pointing \texttt{avrdude} at the \texttt{.hex} you just rebuilt (check
the path/timestamp), not an old one from an earlier exercise. \\
\bottomrule
\end{tabular}
\end{center}

\end{document}
